\clearpage

\section{DESCRIPCIÓN DEL PROBLEMA}

El pensamiento computacional es reconocido internacionalmente como una competencia transversal cuyo desarrollo en etapas escolares contribuye a formar estudiantes con mayor capacidad analítica, lógica y de resolución de problemas. No obstante, 
su incorporación efectiva en el aula requiere más que una declaración de intención pedagógica: exige programas estructurados, recursos didácticos adecuados y mecanismos concretos que articulen a docentes, estudiantes e instituciones educativas escolares en torno a experiencias formativas accesibles y sostenibles.

El Desafío Bebras representa a nivel internacional un modelo probado para atender esta necesidad en el contexto de unidades educativas de nivel primario y secundario. Su enfoque permite que estudiantes de distintos perfiles y edades escolares desarrollen habilidades de pensamiento computacional a través de tareas diseñadas 
sin requerir conocimientos previos de programación. Para que este modelo pueda operar en un país determinado, la organización internacional establece que cada nación debe contar con un organizador nacional responsable de ejecutar el concurso anualmente, mantener una presencia pública del desafío y gestionar los procesos que hacen posible su realización \cite{bebras}. Este modelo no opera mediante una plataforma centralizada compartida, sino a través de sistemas nacionales que cada país desarrolla o adapta según sus propias necesidades organizativas y tecnológicas.

En Bolivia, las condiciones para cumplir con este modelo aún se encuentran en proceso de construcción. Iniciativas previas orientadas a promover el pensamiento computacional entre estudiantes de nivel primario y secundario no lograron consolidarse institucionalmente,
lo que evidencia que el interés por el área no ha sido suficiente por sí solo para sostener programas educativos dirigidos a esta población en colegios y escuelas del país. A esto se suma que Bolivia se incorpora a la red Bebras en 2026 sin la infraestructura organizativa ni tecnológica especializada que caracteriza a los miembros activos de la comunidad internacional.

La exposición marginal de estudiantes de primaria y secundaria al pensamiento computacional en el sistema escolar boliviano se relaciona directamente con la ausencia de programas estructurados que orienten a los docentes en su enseñanza. Sin un marco institucional que articule a los actores del sistema educativo escolar en torno a este propósito, los esfuerzos individuales resultan insuficientes para generar impacto sostenido en niños y jóvenes en edad escolar. Desde el punto de vista operativo, la gestión de un concurso educativo como Bebras requiere procesos coordinados que involucran el registro de instituciones escolares y de estudiantes participantes de nivel primario y secundario, la administración de categorías y tareas por rango etario escolar, la ejecución controlada de pruebas dirigidas a dichos estudiantes y la publicación de resultados. Cuando estos procesos se desarrollan de forma fragmentada o mediante herramientas no diseñadas para este tipo de competencia escolar, la carga operativa aumenta considerablemente y la consistencia de la experiencia para los estudiantes participantes se ve comprometida.

Los efectos de esta situación son visibles tanto en la dimensión educativa como en la organizativa. En el plano educativo, los estudiantes escolares bolivianos cuentan con escasa exposición a experiencias estructuradas de pensamiento computacional, y los docentes disponen de recursos limitados para incorporar estos contenidos en su práctica. El interés de los jóvenes en edad escolar por las ciencias de la computación, aunque presente, no encuentra espacios formales 
que lo orienten y desarrollen. En el plano organizativo, Bolivia se incorpora a la red Bebras sin la base institucional y tecnológica que caracteriza a los países con mayor trayectoria, lo que restringe el alcance de la primera edición nacional 
y condiciona la posibilidad de darle continuidad en años posteriores.

El árbol del problema sintetiza de manera estructurada el problema central, sus causas y los efectos que produce en el contexto educativo boliviano. Su propósito es mostrar que la situación no responde a una dificultad aislada, sino a la 
convergencia de limitaciones organizativas, tecnológicas y educativas que condicionan la viabilidad y sostenibilidad del Desafío Bebras en Bolivia como iniciativa escolar.

\begin{figure}[ht]
\centering
\resizebox{\textwidth}{!}{%
\begin{tikzpicture}[
    font=\small,
    caja/.style={
        draw,
        rectangle,
        rounded corners=4pt,
        align=center,
        minimum height=1cm,
        text width=3.4cm,
        fill=white,
        inner sep=4pt
    },
    cajaGrande/.style={
        draw,
        rectangle,
        rounded corners=4pt,
        align=center,
        minimum height=1.2cm,
        text width=6cm,
        fill=white,
        inner sep=4pt
    },
    cajaCausa/.style={
        draw,
        rectangle,
        rounded corners=4pt,
        align=center,
        minimum height=1cm,
        text width=5cm,
        fill=white,
        inner sep=4pt
    },
    linea/.style={-{Stealth[length=4mm]}, thick}
]

% Efecto superior
\node[cajaGrande] (impacto) at (0,0)
    {Impacto educativo reducido \\del programa};

% Efectos intermedios
\node[caja] (exposicion) at (-9.0,-3.5)
    {Baja exposición\\ estudiantil};
\node[caja] (docentes) at (-4.5,-3.5)
    {Docentes con acceso limitado \\a programas};
\node[caja] (interes) at (0,-3.5)
    {Interés juvenil sin\\ espacios formales};
\node[caja] (alcance) at (4.5,-3.5)
    {Alcance \\organizativo \\restringido};
\node[caja] (continuidad) at (9.0,-3.5)
    {Continuidad reducida\\ del programa};

% Problema central
\node[cajaGrande] (problema) at (0,-7.5)
    {Limitaciones en la \\implementación del Desafío \\Bebras en Bolivia para \\estudiantes escolares};

% Causas intermedias
\node[cajaCausa] (infraestructura) at (-4.5,-11.5)
    {Infraestructura tecnológica\\ no especializada};
\node[cajaCausa] (gestion) at (4.5,-11.5)
    {Gestión fragmentada\\ del concurso};

% Causas raíz
\node[cajaCausa] (presencia) at (-4.5,-14.5)
    {Pensamiento computacional\\ con presencia marginal\\ en el aula};
\node[cajaCausa] (iniciativas) at (4.5,-14.5)
    {Iniciativas previas sin\\ continuidad};

% Flechas efectos intermedios -> efecto superior
\draw[linea] (exposicion.north) -- (impacto.south west);
\draw[linea] (docentes.north) -- (impacto.south west);
\draw[linea] (interes.north) -- (impacto.south);
\draw[linea] (alcance.north) -- (impacto.south east);
\draw[linea] (continuidad.north) -- (impacto.south east);

% Flechas problema central -> efectos intermedios (CORREGIDO)
\draw[linea] (problema.north west) -- (exposicion.south);
\draw[linea] (problema.north west) -- (docentes.south);
\draw[linea] (problema.north) -- (interes.south);
\draw[linea] (problema.north east) -- (alcance.south);
\draw[linea] (problema.north east) -- (continuidad.south);

% Flechas causas intermedias -> problema central
\draw[linea] (infraestructura.north) -- (problema.south west);
\draw[linea] (gestion.north) -- (problema.south east);

% Flechas causas raíz -> causas intermedias
\draw[linea] (presencia.north) -- (infraestructura.south);
\draw[linea] (iniciativas.north) -- (gestion.south);

\end{tikzpicture}%
}
\caption{Árbol del problema para la implementación del Desafío Bebras Bolivia}
\label{fig:arbol-problemas}
\end{figure}

\subsection{DEFINICIÓN DEL PROBLEMA}

El problema central radica en las limitaciones en la implementación del Desafío Bebras en Bolivia para estudiantes de nivel primario y secundario, en un contexto donde factores educativos, organizativos y tecnológicos condicionan su desarrollo y continuidad entre los niños y jóvenes en edad escolar del país.